\section{Experimental protocol}
\label{sec:protocol}

The computational protocol (partitions and seeds, base models, design, surfaces and gate thresholds, optimizers and
scoring) was committed to version control before the replicated study was launched with the planned ten
replications per dataset. Its only change during execution was a solver-robustness fix after a crash in
replication~0; the replication-0 stages completed before the fix were kept, and the part of the change that can be
checked on the stored surfaces moves no stored solution by more than $10^{-10}$ in any weight (supplementary Section~S1). After the ten-replication
analysis had been completed and reported, the study was extended to 30 replications, reusing the computational
protocol without modification, and \cref{sec:results:r30} reports what that extension changed. The NSGA-II baseline
(\cref{sec:results:nsga2}) was added after the 30-replication analysis, with its configuration committed before its
first run.

\subsection{Datasets}
\label{sec:protocol:data}

Four public binary tabular classification problems span prevalence, dimensionality and feature type
(\cref{tab:datasets}): Santander \citep{kagglesantander} (high-dimensional, anonymized, numeric, 10\% positives);
BNP Paribas \citep{kagglebnp} (numeric plus high-cardinality categorical features, 76\% positive); Porto Seguro
\citep{kaggleporto} (strongly imbalanced at 3.6\% positives, subsampled to 200{,}000 rows by a single stratified
draw with a recorded seed to bound runtime); and UCI credit default \citep{yeh2009ucidataset} (small and
low-dimensional). For the three Kaggle competition datasets the competition test labels are not
public and were never used: all training, weighting, thresholding and model selection use only the published training
file, from which every performance figure in this paper is derived.

\begin{table}[t]
\centering
\caption{Datasets. Rows used, number of raw features, positive-class prevalence and training/holdout sizes of each outer partition (stratified 80/20). Kaggle test labels are never used; all selection uses 5-fold out-of-fold predictions within the training part. $^{\dagger}$200,000-row stratified subsample of the 595,212 available rows, drawn once with a recorded seed.}
\label{tab:datasets}
\small
\resizebox{\textwidth}{!}{%
\begin{tabular}{llrrrl}
\toprule
Dataset & Source & Rows & Features & Prevalence & Train / holdout rows \\
\midrule
Santander & Santander Customer Transaction (Kaggle 2019) & 200,000 & 200 & 0.100 & 160,000 / 40,000 \\
BNP Paribas & BNP Paribas Cardif Claims (Kaggle 2016) & 114,321 & 131 & 0.761 & 91,456 / 22,865 \\
Porto Seguro & Porto Seguro Safe Driver (Kaggle 2017)$^{\dagger}$ & 200,000 & 57 & 0.036 & 160,000 / 40,000 \\
UCI credit & Default of credit card clients (UCI id 350) & 30,000 & 23 & 0.221 & 24,000 / 6,000 \\
\bottomrule
\end{tabular}}
\end{table}


\subsection{Base models and leakage control}
\label{sec:protocol:models}

The ensemble contains $M = 5$ deliberately heterogeneous learners with fixed hyperparameters and no hyperparameter
search: $\ell_2$-regularized logistic regression; Gaussian naive Bayes; $k$-nearest neighbours with $k = 100$; a
random forest with 200 trees and a minimum of 20 samples per leaf; and gradient-boosted trees (XGBoost, 400 rounds,
depth 6, learning rate 0.1, subsample and column subsample 0.8) \citep{chen2016xgboost}. The first four are the
\texttt{scikit-learn} implementations \citep{pedregosa2011scikit}. These values were committed to the project code on
2026-07-06, before the replicated study, and were not changed afterwards; Gaussian naive Bayes uses the library
defaults, as do the inverse regularization strength of logistic regression ($C = 1$) and the XGBoost depth. The zoo
spans the accuracy--cost--calibration space rather than forming a competitive stack: Gaussian naive Bayes is the best
member out of fold on Santander but has the worst log-loss on the other three datasets, and $k$-nearest neighbours,
the weakest member by ROC-AUC on Santander and Porto Seguro, costs more at inference time than the other four
combined (supplementary Section~S2; implications in \cref{sec:limitations}). Preprocessing (imputation,
standardization and encoding, detailed per model in supplementary Section~S1) is fitted inside each training fold and
never on validation data.

Each replication splits the dataset by a stratified 80/20 partition into a training part and an untouched holdout.
Within the training part, stratified 5-fold cross-validation produces out-of-fold probability estimates $\mathbf{P}$
for all five models; these are the only quantities the design, the surfaces, the optimizers and the selection rules
ever see. Each model is then refitted once on the whole training part to produce holdout probabilities $\mathbf{Q}$,
used exclusively for the confirmation analysis of \cref{sec:results:holdout}. Every entry of $\mathbf{P}$ and
$\mathbf{Q}$ is clipped to $[10^{-3}, 1-10^{-3}]$, a constant fixed before the study and never changed: $k$-nearest
neighbours with $k = 100$ outputs exact zeros and ones, which can make the log-loss infinite at its vertex, and
clipping the components rather than the blend keeps \cref{eq:ensemble} linear in $\wvec$. Because \cref{eq:ensemble} is linear in
$\wvec$, evaluating any weight vector costs one matrix--vector product, which is what makes an empirical reference of
$10^{5}$ points and a metamodel-free NBI arm affordable.

Inference cost $c_i$ is measured, not assumed: per model and replication, the median wall-clock time of five
consecutive \texttt{predict\_proba} calls of the refitted model, including preprocessing, in milliseconds per
1{,}000 rows. One batch of holdout rows, drawn with the replication seed, is shared by the five models (10{,}000 rows;
the whole 6{,}000-row holdout on UCI credit). The models are timed one after another in a fixed, non-randomized order
(logistic regression, Gaussian naive Bayes, $k$-nearest neighbours, random forest, XGBoost), each right after its own
refit, an untimed holdout prediction and an untimed warm-up call on the batch. Three of them ($k$-nearest neighbours,
the random forest and XGBoost) are set to use all 14 cores (\texttt{n\_jobs}~$=-1$); the code and the launch script
set no other parallel option or thread limit for the timed calls.

\subsection{Replication, seeds and the statistical unit}
\label{sec:protocol:replication}

The entire pipeline, from the outer partition through scoring, is executed $R = 30$ times per dataset, giving 120
replications, with the outer partition, the inner folds and every stochastic component seeded by
$\text{20260904} + r$ for replication $r$. All seeds are persisted in the manifest, and the checkpointed runner
described in the Reproducibility statement recomputes any stage exactly, except that re-running the out-of-fold stage
re-measures the wall-clock costs $c_i$ and therefore changes every result that depends on them.

The statistical unit is \textbf{one outer partition of one dataset}. Because all methods are evaluated on the same
partition, every comparison is \emph{paired}. Two consequences constrain the analysis throughout.

First, the 30 partitions of a dataset overlap heavily: two independent stratified 80/20 draws have an expected
intersection of $0.8 \times 0.8 = 64\%$ of the dataset, which is $80\%$ of each training part, a Jaccard index of
two-thirds. Any two training parts therefore share most of their rows.
The replication distribution therefore measures \emph{sensitivity to data partitioning}, not sampling from a
population, and the effective information is far below 30 independent runs. Second, and consequently, the dataset is
the unit of generalization: results are summarized per dataset and the four dataset-level results are then compared.
Nothing is pooled across the $4 \times 30$ replications, and no statement of the form ``$N = 120$ independent
experiments'' is made anywhere in this paper.

\subsection{Statistical analysis}
\label{sec:protocol:stats}

The statistical analysis plan (supplementary Section~S1) designates four primary comparisons. NBI-B versus NBI-A and
NBI-C versus NBI-B address RQ4 and RQ5. NBI-C versus random scalarization and NBI-C versus the random Dirichlet sample
compare NBI-C, as specified and not at equal budget, with the two cheap alternatives of \cref{sec:method:refs}: the
first comparison is confounded with the objective source and the second is a floor check. The plan was committed to
version control after the first ten replications had been run and analysed, minutes after the extension to 30
replications was launched and before any of its partitions had been completed, so it is prospective only for
partitions 10--29. The paired contrasts coded before the benchmark produced its first output were descriptive and
compared both cheap alternatives with NBI-B rather than NBI-C; \cref{sec:results:metamodelfree} reports both
pairings. The
primary endpoints are $\igdp$ and the hypervolume ratio, computed on the revalidated \emph{out-of-fold objectives of
the 80\% training part} (the only quantities the design, the surfaces, the optimizers and the selection rules ever
see; \cref{sec:protocol:models}) against the empirical reference of the same replication, under the weighted cost
that the optimizers actually minimized; the support-cost repetition is a sensitivity analysis. The
untouched 20\% holdout enters no primary endpoint: it is used only for the confirmation analysis of
\cref{sec:results:holdout}. The primary endpoints therefore measure how well a returned front approximates the
empirical front on the training part's out-of-fold objectives, not how a front generalizes; only the holdout
transfer bears on generalization, and it scores one selected solution per set. All other indicators are descriptive.

Paired differences are signed so that $\Delta > 0$ always favours the first-named (new) method:
$\Delta\hv = \hv_{\text{new}} - \hv_{\text{ref}}$ and
$\Delta\igdp = \igdp_{\text{ref}} - \igdp_{\text{new}}$. For each dataset $\times$ comparison $\times$ endpoint we
report the mean and median with percentile-bootstrap 95\% intervals (10{,}000 resamples, fixed seed), the quartiles
and range \citep{efron1993bootstrap}, the win/tie/loss counts with a tie tolerance of $10^{-4}$, the win fraction
with a Wilson 95\% interval \citep{wilson1927probable}, and the matched-pairs rank-biserial correlation
\citep{kerby2014simple} as a distribution-free effect size. These and all other replication-level intervals in this
paper treat the replications of a dataset as independent observations: they describe partition sensitivity within
that dataset, carry no correction for partition overlap and do not by themselves establish significance.

An ordinary paired $t$-test on resampled partitions is anticonservative, so the primary test is the Nadeau--Bengio
corrected resampled $t$-test \citep{nadeau2003inference}: with $J = 30$ paired differences $d_j$,
\begin{equation}
t \;=\; \frac{\bar d}{\sqrt{\operatorname{var}(d)\left(\tfrac{1}{J} + \rho\right)}}, \qquad
\rho \;=\; \frac{n_{\text{test}}}{n_{\text{train}}} \;=\; \frac{0.20}{0.80} \;=\; 0.25,
\label{eq:nb}
\end{equation}
on $J - 1$ degrees of freedom. The correction inflates the variance by a factor of
$(1/30 + 0.25)/(1/30) \approx 8.5$, multiplying standard errors by about 2.9 relative to the naive test. One caveat
attaches to $\rho$ and we state it plainly: the heuristic $\rho = n_{\text{test}}/n_{\text{train}}$ was derived for a
statistic measured on the held-out portion of each resample, whereas both primary endpoints here are measured on the
training part, whose rows two replications share to the extent computed above. The value $\rho = 0.25$ is therefore
an assumed correction, not an estimate of the correlation between replications, which was not measured; the
training-part overlap suggests that it may understate that correlation, in which case the corrected test would
still be anticonservative. Every significance statement below holds \emph{at} the assumed $\rho = 0.25$, and the
effect sizes, the paired direction counts and their spread over partitions, which do not depend on $\rho$, are the
primary evidence and are reported alongside it. Within each
dataset the eight primary tests form one family and receive a Holm correction \citep{holm1979simple}; no correction
is applied across datasets, because datasets are not pooled. A Wilcoxon signed-rank test
\citep{wilcoxon1945individual} is reported alongside as a distribution-free check, but it is not corrected for overlap and cannot rescue a non-significant corrected result.

Where a paired-difference distribution is visibly bimodal or heavy-tailed, the mean-based test is uninformative by
construction and we say so, reporting the median with its bootstrap interval, the win fraction with its Wilson
interval and the rank-biserial correlation as descriptive evidence of consistency instead. Throughout, effect consistency takes precedence
over $p$-values. Frequencies over replications are reported with Wilson and Jeffreys 95\% intervals \citep{brown2001interval}.
